\chapter{Pore analysis}

%\section{Pore analysis using \acs{LTDSC}, \acs{MIP}, \acs{CT} and \acs{SEM}}
The determination of porosity is fundamental to understanding and predicting the performance and durability of concrete.
Whilst macroscopic air voids influence properties like freeze-thaw resistance, the gel porosity dictates the material's intrinsic characteristics.
For example, the transport mechanisms (permeability and diffusion) of water and ions are determined by the pore network caused by the fine structures inherent of the \ac{CSH} phase.
Accurately quantifying this gel porosity is essential to assess the resistance to chemical attack (e.g. chloride ingress or sulphate attack), and its mechanical performance.

On a macroscopic scale, comparing different pore measurement methods is feasible.
However, as the pore size decreases, this comparison becomes increasingly challenging due to physical limits and the fractal nature of porosity.
To address this, a standard material was utilised to benchmark and compare selected methods.

\section{Porous glass standard}

Two porous glass standards (\textsc{VYCOR} glass with \num{5} and \qty{20}{\nano\meter} diameter, see \zcref{sec:boraglas}) were evaluated using \ac{SEM}, \ac{LTDSC}, and \ac{HNMR} spectroscopy.
The glass standards were initially defined using \ac{MIP} by the manufacturer (\zcref{fig:vycor-mip}). % or \ce{N2}-ad/absorption.
% TODO: die Endporosität wird normalerweise eher durch Engstellen (necks) definiert als durch die Bulk-porosität

\subsection{\acs*{SEM}}

To obtain \ac{SEM} images, the glass was fractured and even surfaces were imaged without any coating.
As the \ac{SEM} images indicate (\zcref{fig:ltdsc-vycor-glass-a, fig:ltdsc-vycor-glass-b}), the pores are interconnected, which is plausible due to its intended use as a filter material and the \textsc{VYCOR} manufacturing process used (sintering of glass powder).

\begin{figure}[htb!]
    \centering

	\begin{subfigure}{0.49\textwidth}
		\includegraphics[width=\textwidth]{ltdsc/glass/segmentation/5nm.pdf}
		\caption{Raw, pore size \qty{5}{\nano\meter} (diameter)}
		\label{fig:ltdsc-vycor-glass-a}
	\end{subfigure}
	\hfill
	\begin{subfigure}{0.49\textwidth}
		\includegraphics[width=\textwidth]{ltdsc/glass/segmentation/20nm.pdf}
		\caption{Raw, pore size \qty{20}{\nano\meter} (diameter)}
		\label{fig:ltdsc-vycor-glass-b}
	\end{subfigure}

	\begin{subfigure}{0.49\textwidth}
		\includegraphics[width=\textwidth]{ltdsc/glass/20nm_005old_otsu.png}
		\caption{Segmented, pore size \qty{5}{\nano\meter} (diameter)}
		\label{fig:ltdsc-vycor-glass-c}
	\end{subfigure}
	\hfill
	\begin{subfigure}{0.49\textwidth}
		\includegraphics[width=\textwidth]{ltdsc/glass/5nm_004old_otsu.png}
		\caption{Segmented, pore size \qty{20}{\nano\meter} (diameter)}
		\label{fig:ltdsc-vycor-glass-d}
	\end{subfigure}


    % \begin{subfigure}{.495\textwidth}
    %   \centering
    %   \begin{tikzpicture}
    %     \begin{semilogxaxis}[
    %             legend style={
    %                 at={(0.05,.8)},
    %                 anchor=west,
    %                 cells={align=left}
    %             },
    %             xmin=.64, xmax=160,
    %             ymin=0, ymax=1.1,
    %             width=1.1\linewidth,
    %             height=6cm,
    %             grid=major,
    %             xlabel={pore diameter in \unit{\nano\meter}},
    %             ylabel={normalised frequency},
    %             grid style={line width=.1pt, draw=gray!10},
    %             yticklabel style={
    %                     /pgf/number format/fixed,
    %                     /pgf/number format/precision=3
    %             },
    %             xticklabels={2.5,5,10,20,40,80,160,320},
    %             xtick={1.25, 2.5,5,10,20,40,80,160},
    %             ymajorticks=false,
    %             scaled y ticks=false,
    %             scaled x ticks=false,
    %         ]

    %         % --- ANGEPASSTER PLOT ---
    %         \addplot [
	% 			no marks,
    %             ybar,
    %             fill=darkgreen,
    %             draw=darkgreen,
    %             bar width=10pt,
    %             bar shift=-5pt,
    %             area legend
    %         ] table [x=x, y=y_norm] {
    %             x      y       y_norm
    %             2.5    0       0
    %             5      536     0.2902
    %             10     1847    1
    %             20     1502    0.8134
    %             40     51      0.0276
    %             80     0       0
    %         };
	% 		\addlegendentry{image \num{1}};

    %         % --- ANGEPASSTER PLOT (statt dashed) ---
    %         \addplot [
	% 			no marks,
    %             ybar,
    %             pattern=north east lines,
    %             pattern color=darkgreen,
    %             draw=darkgreen,
    %             bar width=10pt,
    %             bar shift=5pt,
    %             area legend
    %         ] table [x=x, y=y_norm] {
    %             x      y       y_norm
    %             0.64   0       0
    %             1.25   3       0.0045
    %             2.5    132     0.1959
    %             5      110     0.1633
    %             10     442     0.6567
    %             20     674     1
    %             40     90      0.1336
    %             80     1       0.0015
    %             160    0       0
    %         };
	% 		\addlegendentry{image \num{2}};
    %     \end{semilogxaxis}
    %   \end{tikzpicture}
    %   \caption{
	% 	Pore size distribution \qty{5}{\nano\meter} porous \textsc{VYCOR}-glass.
	%   }
	%   \label{fig:ltdsc-vycor-glass-e}
    % \end{subfigure}
    % \hfill
    % \begin{subfigure}{.495\textwidth}
    %     \centering
    %     \begin{tikzpicture}
    %       \begin{semilogxaxis}[
    %             legend style={
    %                 at={(0.05,.8)},
    %                 anchor=west,
    %                 cells={align=left}
    %             },
    %             xmin=.64, xmax=160,
    %             ymin=0, ymax=1.1,
    %             width=1.1\linewidth,
    %             height=6cm,
    %             grid=major,
    %             xlabel={pore diameter in \unit{\nano\meter}},
    %             ylabel={normalised frequency},
    %             grid style={line width=.1pt, draw=gray!10},
    %             yticklabel style={
    %                     /pgf/number format/fixed,
    %                     /pgf/number format/precision=3
    %             },
    %             xticklabels={2.5,5,10,20,40,80,160,320},
    %             xtick={1.25, 2.5,5,10,20,40,80,160},
    %             ymajorticks=false,
    %             scaled y ticks=false,
    %             scaled x ticks=false,
    %             yticklabel pos=right,
    %         ]
    %         % --- ANGEPASSTER PLOT ---
    %         \addplot+ [
	% 			no marks,
    %             ybar,
    %             fill=orange,
    %             draw=orange,
    %             bar width=10pt,
    %             bar shift=-5pt,
    %             area legend
    %         ] table [x=x, y=y_norm] {
    %             x      y       y_norm
    %             1.25   0       0
    %             2.5    35      0.0429
    %             5      116     0.1422
    %             10     349     0.4272
    %             20     816     1
    %             40     136     0.1667
    %             80     0       0
    %         };
	% 		\addlegendentry{image \num{1}};

    %         % --- ANGEPASSTER PLOT (statt dashed) ---
    %         \addplot+ [
	% 			no marks,
    %             ybar,
    %             pattern=north east lines,
    %             pattern color=orange,
    %             draw=orange,
    %             bar width=10pt,
    %             bar shift=5pt,
    %             area legend
    %         ] table [x=x, y=y_norm] {
    %             x      y       y_norm
    %             1.25   0       0
    %             2.5    1       0.0004
    %             5      314     0.1180
    %             10     1244    0.4675
    %             20     2661    1
    %             40     336     0.1263
    %             80     0       0
    %         };
	% 		\addlegendentry{image \num{2}};


    %       \end{semilogxaxis}
    %     \end{tikzpicture}
    %     \caption{
	% 		Pore size distribution of \qty{20}{\nano\meter} porous \textsc{VYCOR} glass.
	% 	}
	% 	\label{fig:ltdsc-vycor-glass-f}
    % \end{subfigure}

 	\begin{subfigure}{.495\textwidth}
      \begin{tikzpicture}
        \begin{semilogxaxis}[
				legend style={
					at={(0.985,.8)},
					anchor=east,
					cells={align=left},
					font=\footnotesize
				},
				xmin=3, xmax=41,
				ymin=0, ymax=1,
				width=\linewidth,
				height=6cm,
				grid=major,
				xlabel={pore diameter $d_\text{P}$ in \unit{\nano\meter}},
				ylabel={normalised frequency},
				grid style={line width=.1pt, draw=gray!10},
				yticklabel style={
						/pgf/number format/fixed,
						/pgf/number format/precision=3
				},
				xticklabels={2.5,5,10,20,40},
				xtick={2.5,5,10,20,40},
				ytick={0,-0.1,-0.2},
				scaled y ticks=false,
				scaled x ticks=false,
            ]

			  \addplot [ no marks, orange ] table [x index=0, y index=1, col sep=comma ] {figures/ltdsc/glass/20nm_005old_otsu.csv};
			  \addlegendentry{specimen 1};

        \end{semilogxaxis}
	  \end{tikzpicture}
	  \caption{Pore size distribution of \qty{5}{\nano\meter} porous \textsc{VYCOR} glass.}
	  \label{fig:ltdsc-vycor-glass-e}

	\end{subfigure}
	\hfill
	\begin{subfigure}{.495\textwidth}

		\begin{tikzpicture}
		  \begin{semilogxaxis}[
				legend style={
					at={(0.015,.8)},
					anchor=west,
					cells={align=left},
					font=\footnotesize
				},
				xmin=3, xmax=41,
				ymin=0, ymax=1,
				width=\linewidth,
				height=6cm,
				grid=major,
				xlabel={pore diameter $d_\text{P}$ in \unit{\nano\meter}},
				ylabel={normalised frequency},
				grid style={line width=.1pt, draw=gray!10},
				yticklabel style={
						/pgf/number format/fixed,
						/pgf/number format/precision=3
				},
				xticklabels={2.5,5,10,20,40},
				xtick={2.5,5,10,20,40},
				ytick={0,-0.1,-0.2},
				scaled y ticks=false,
				scaled x ticks=false,
				yticklabel pos=right,
			]

			\addplot [ no marks, darkgreen] table [x index=0, y index=1, col sep=comma] {figures/ltdsc/glass/5nm_004old_otsu.csv};
			\addlegendentry{specimen 1};


		  \end{semilogxaxis}
		\end{tikzpicture}
		\caption{Pore size distribution of \qty{20}{\nano\meter} porous \textsc{VYCOR} glass.}
	  	\label{fig:ltdsc-vycor-glass-f}

	\end{subfigure}
    \caption{
		\ac{SEM}-\ac{SE} images of two \textsc{VYCOR}-glass samples, their respective segmentation result and the pore size distribution.
		The pore size was determined using two segmented \ac{SEM} image for each glass type.
	} % Probenbilder-rodaten vmtl vertauscht.
    \label{fig:ltdsc-vycor-glass}
\end{figure}


% C:\Users\Florian Kleiner\Nextcloud\Uni\WdB\Paper\2024 LT-DSC-XRD-BSE Alit-Belit Paper\5. REM+EDX\Nanoporöses Glas\2022_02_08 Poröses Glas

Since the dwell-time has to be chosen as short as possible to avoid charging, the images became very noisy and were denoised using a median blur filter using a \numproduct{3x3} kernel.

The pores can be segmented using thresholding to determine the pore size distribution (see \zcref{fig:ltdsc-vycor-glass-c,fig:ltdsc-vycor-glass-d}). %  (see \zcref{sec:segmentation})

Areas with obvious damages due to the fracture process were excluded from the analysis.
The resulting pore size distributions of two images per specimen are shown in \zcref{fig:ltdsc-vycor-glass-e, fig:ltdsc-vycor-glass-f}.
As indicated by the inhomogeneous distribution of the segmented pores (\zcref{fig:ltdsc-vycor-glass-c}), the selected threshold value is not producing an uniform result.

Furthermore, the measured pore radii measured using this traditional thresholding method are off by a factor of \num{4} and more (\qty{5}{\nano\meter} glass) and \num{2} (\qty{20}{\nano\meter} glass) compared to the \ac{MIP} measurements provided by the manufacturer (\zcref{fig:vycor-mip}).

% TODO: prüfe, ob tatsächlich Porenradien gemessen wurden!


\begin{figure}[htbp]
	\centering
	\begin{subfigure}[b]{0.46\textwidth}

		\begin{tikzpicture}
			\begin{axis}[
				xlabel={Median \textsc{Feret} diameter in \unit{\nano\meter}},
				ylabel={blocksize in \unit{\nano\meter}},
				colormap name=buw, 		 % Choose a colormap (others: hot, cool, jet, etc.)
				colorbar=false,                       % Display the color bar
				view={0}{90},                   % Set view angles for a 2D top-down view
				enlargelimits=false,            % Fit plot tightly to data range
				axis on top,
				point meta min=0,
				xmin=0, xmax=40,
				ymin=0, ymax=190,
				width=7.2cm,
				height=6cm,
			]
				\addplot [
					matrix plot*,         % Use matrix plot for heatmaps (starred version often better for coordinate input)
					point meta=explicit,  % Color data is explicitly given in the 'meta' column
					shader=flat,
					mesh/cols=49,% 50, %118
				]
				table [
					x={Feret diameter [nm]}, % Column name for X coordinate
					y={Blocksize (nm)},      % Column name for Y coordinate
					meta={Normalized Frequency},        % Column name for the color value
					col sep=comma,           % Specify comma as the column separator
					header=true              % Indicate that the CSV file has a header row
				] {figures/mip/vycor5nm.crop.csv}; % Path to your CSV file
			\end{axis}
		\end{tikzpicture}
		\caption{\qty{5}{\nano\meter} porous \textsc{VYCOR} glass.}
		\label{fig:mip_heatmap_vycor5nm}

	\end{subfigure}
	\hfill
	\begin{subfigure}[b]{0.52\textwidth}

		\begin{tikzpicture}
			\begin{axis}[
				xlabel={Median \textsc{Feret} diameter in \unit{\nano\meter}},
				%ylabel={blocksize in \unit{\nano\meter}},
				colorbar,                       % Display the color bar
				colorbar style={
					ylabel={normalised frequency},
					ytick=\empty,
				}, % Label for the color bar
				colormap name=buw, 		 % Choose a colormap (others: hot, cool, jet, etc.)
				view={0}{90},                   % Set view angles for a 2D top-down view
				enlargelimits=false,            % Fit plot tightly to data range
				point meta min=0,               % Optional: Set minimum value for the color axis
				axis on top,
				xmin=0, xmax=40,
				ymin=0, ymax=190,
				yticklabel pos=right,
				width=7.2cm,
				height=6cm,
				ytick=\empty,
			]
				\addplot [
					matrix plot*,
					point meta=explicit,
					shader=flat,
					mesh/cols=49,
				]
				table [
					x={Feret diameter [nm]},
					y={Blocksize (nm)},
					meta={Normalized Frequency},
					col sep=comma,
					header=true
				] {figures/mip/vycor20nm.crop.csv};
			\end{axis}
		\end{tikzpicture}
		\caption{\qty{20}{\nano\meter} porous \textsc{VYCOR} glass.}
		\label{fig:mheatmap_vycor20nm}

	\end{subfigure}
	\caption{
		Heatmaps indicating the influence of the `blocksize' option using `adaptive threshold' on the measured median \textsc{Feret} diameter of the \textsc{VYCOR} glass.
		The frequency is normalised per blocksize.
		The blocksize is the size of the sliding window used for the adaptive thresholding.
	}
	\label{fig:heatmap_vycor}
\end{figure}


Since \ac{SEM}-\ac{SE} images tend to have brightness variations from one side to the other e.g. due to slightly tilted surfaces, adaptive thresholding was applied.
Adaptive thresholding is a segmentation method that adjusts the threshold value dynamically based on the local properties of the image, such as the intensity within a sliding window.
The size of this window is a variable called blocksize which has to be set manually.
For each segmented pore, the median \textsc{Feret} diameter was calculated.
The resulting median \textsc{Feret} diameter distribution was then normalised to the maximum value within the distribution.
Depending on the blocksize used, very different results were obtained, as shown in \zcref{fig:heatmap_vycor}.

For \textsc{VYCOR}-glass with a given porosity of \qty{20}{\nano\meter}, using a blocksize larger than \qty{50}{\nano\meter} produced reliable results, with the main median \textsc{Feret} diameter around \qty{20}{\nano\meter}.
However, for \qty{5}{\nano\meter} glass, the porosity was overestimated by a factor of \num{2.5}, with a median \textsc{Feret} diameter around \qty{15}{\nano\meter} (see \zcref{fig:mip_heatmap_vycor5nm}).
For block sizes smaller than \qty{50}{\nano\meter}, the segmentation result was very noisy, and the \textsc{Feret} diameter was underestimated and unreliable.
It is recommended to use a blocksize that is at least three times larger than the expected feature size to achieve accurate segmentation results.

% TODO: Median pore size eintragen ins Diagram
% TODO: normales thresholding auf eines der identischen Bilder anwenden zum Vergleich.
% TODO: Denoising erwähnen
% TODO: neues Bild mit Auswirkung von Blocksize



\FloatBarrier
\subsection{\acs*{LTDSC}}

The glass standards were measured using a simplified temperature regimen.
It also experienced ice crystal-inducing freeze-thaw cycles before the actual measurement.
For the measurement, the water saturated specimen was then cooled down to \qty{-60}{\celsius} and heated up to \qty{20}{\celsius} at a cooling and heating rate of $\pm$\,\qty{2}{\kelvin\per\minute}.

\begin{figure}[!htb]
    \centering

 	\begin{subfigure}{.48\textwidth}
      \begin{tikzpicture}
        \begin{semilogxaxis}[
				blue,
				axis line style={black},
				xmin=1.5, xmax=60,
				ymin=-.20, ymax=0,
				y dir=reverse,
				width=\linewidth,
				height=6cm,
				grid=none,
				xlabel={temperature in \unit{\celsius}},
				xticklabels={\num{-51.2}, \num{-25.3}, \num{-12.4}, \num{-4.6}, \num{-2.0}},
				xtick={2.5,5,10,25,50},
				ytick=\empty,
				scaled y ticks=false,
				scaled x ticks=false,,
				axis x line*=top
            ]

        \end{semilogxaxis}
        \begin{semilogxaxis}[
				legend style={
					at={(0.015,.8)},
					anchor=west,
					cells={align=left},
					font=\footnotesize
				},
				xmin=1.5, xmax=60,
				ymin=-.20, ymax=0,
				y dir=reverse,
				width=\linewidth,
				height=6cm,
				grid=major,
				xlabel={pore radius $r_\text{P}$ in \unit{\nano\meter}},
				ylabel={heat flow $\dot{Q}$ in \unit{\joule\per\second\per\gram}},
				grid style={line width=.1pt, draw=gray!10},
				yticklabel style={
						/pgf/number format/fixed,
						/pgf/number format/precision=3
				},
				xticklabels={2.5,5,10,25,50},
				xtick={2.5,5,10,25,50},
				ytick={0,-0.1,-0.2},
				scaled y ticks=false,
				scaled x ticks=false,
            ]
			  \addplot[dotted, thick, samples=2, darkgrey,forget plot] coordinates {(3.19,-.2)(3.19,.2)};
			  \node[darkgrey, rotate=90, anchor=west] at (axis cs: 2.6,-.06) {\scriptsize <\,\qty{-40}{\celsius}!};

			\addplot [ no marks, darkgreen] table [each nth point=4, filter discard warning=false, unbounded coords=discard,x index=2, y index=1, col sep=comma] {figures/ltdsc/glass/ExpDat_2.1-5nm_freeze_segment_18.csv};
			\addlegendentry{specimen 1};
			\addplot [ no marks, darkgreen, dashed] table [each nth point=4, filter discard warning=false, unbounded coords=discard,x index=2, y index=1, col sep=comma] {figures/ltdsc/glass/ExpDat_2.2-5nm_freeze_segment_18.csv};
			\addlegendentry{specimen 2};

        \end{semilogxaxis}
	  \end{tikzpicture}
	  \caption{\qty{5}{\nano\meter} porous \textsc{VYCOR} glass.}

	\end{subfigure}
	\hfill
	\begin{subfigure}{.48\textwidth}

		\begin{tikzpicture}
			\begin{semilogxaxis}[
				blue,
				axis line style={black},
					xmin=1.5, xmax=60,
					ymin=-.20, ymax=0,
					y dir=reverse,
					width=\linewidth,
					height=6cm,
					grid=none,
					xlabel={temperature in \unit{\celsius}},
					xticklabels={\num{-51.2}, \num{-25.3}, \num{-12.4}, \num{-4.6}, \num{-2.0}},
					xtick={2.5,5,10,25,50},
					ytick=\empty,
					scaled y ticks=false,
					scaled x ticks=false,
					axis x line*=top,
				]

		  \end{semilogxaxis}
		  \begin{semilogxaxis}[
				legend style={
					at={(0.015,.8)},
					anchor=west,
					cells={align=left},
					font=\footnotesize
				},
				xmin=1.5, xmax=60,
				ymin=-.20, ymax=0,
				y dir=reverse,
				width=\linewidth,
				height=6cm,
				grid=major,
				xlabel={pore radius $r_\text{P}$ in \unit{\nano\meter}},
				ylabel={heat flow $\dot{Q}$ in \unit{\joule\per\second\per\gram}},
				grid style={line width=.1pt, draw=gray!10},
				yticklabel style={
						/pgf/number format/fixed,
						/pgf/number format/precision=3
				},
				xticklabels={2.5,5,10,25,50},
				xtick={2.5,5,10,25,50},
				ytick={0,-0.1,-0.2},
				scaled y ticks=false,
				scaled x ticks=false,
				yticklabel pos=right,
			]
			  \addplot[dotted, thick, samples=2, darkgrey,forget plot] coordinates {(3.19,-.2)(3.19,0)};
			  \node[darkgrey, rotate=90, anchor=west] at (axis cs: 2.6,-.06) {\scriptsize <\,\qty{-40}{\celsius}!};

			  \addplot [ no marks, orange ] table [ each nth point=4, filter discard warning=false, unbounded coords=discard, x index=2, y index=1, col sep=comma ] {figures/ltdsc/glass/ExpDat_1.1-20nm_freeze_segment_18.csv};
			  \addlegendentry{specimen 1};
			  \addplot [ no marks, orange, dashed ] table [ each nth point=4, filter discard warning=false, unbounded coords=discard, x index=2, y index=1, col sep=comma ] {figures/ltdsc/glass/ExpDat_1.2-20nm_freeze_segment_18.csv};
			  \addlegendentry{specimen 2};

		  \end{semilogxaxis}
		\end{tikzpicture}
		\caption{\qty{20}{\nano\meter} porous \textsc{VYCOR} glass.}

	\end{subfigure}
    \caption{
		\ac{LTDSC}-measurements of two nano-porous glass types.
		Each type was measured using two specimens. Note that pore sizes below \qty{3.19}{\nano\meter} are undefined as noted in \zcref{eq:brun_freeze}.
	}
    \label{fig:ltdsc_result_glass_freeze}
\end{figure}


This resulted in the following freeze curves as plotted in \zcref{fig:ltdsc_result_glass_freeze}.
The glass type specified with a pore diameter of \qty{5}{\nano\meter} has two peaks in the heat-flow signal.
According to \zcref{eq:brun_freeze} \cite{brun.1977}, a heat flow signal was measured in the pore radius range of \qtyrange{3}{8}{\nano\meter} (pore diameter of \num{6} and \qty{16}{\nano\meter}) with peaks at \num{4} and \qty{6.5}{\nano\meter} (pore diameter of \num{8} and \qty{13}{\nano\meter}).

A heat flow signal was observed for the glass with a specified pore radius of \qty{10}{\nano\meter} for temperatures which can be converted to radii of \num{5} and \qty{12}{\nano\meter}, with peaks at pore radii of \num{7} and \qty{10}{\nano\meter}.

These results show, that the pore diameter is overestimated by a factor of up to \num{2} for the \qty{5}{\nano\meter} glass compared to the \ac{MIP} results (see \zcref{fig:vycor-mip}), while the \qty{20}{\nano\meter} glass shows a good fit for its main peak.
Despite this deviation, the difference in the pore size is still observable.
Compared to the \ac{SEM} results (see \zcref{fig:vicor-SEM}), the \ac{LTDSC} results fit better to the \ac{MIP} results.
Similar to the \ac{SEM} data, the \qty{5}{\nano\meter} glass also exhibits a shift towards larger pore sizes compared to the \ac{MIP} results.


\subsection{\acs*{HNMR}}

%TODO: -> Franz + naber!


\FloatBarrier
\section{Porosity of hydrated \acs*{C2S} and \acs*{C3S}}

In contrast to porous glass, the generation of fracture surfaces in hydrated \ac{C2S} and \ac{C3S} is not sufficient for a pore analysis using \ac{SEM}-\ac{BSE} images.
Therefore, a good specimen preparation is important  to obtain a usable pore analysis.

\subsection{\acs*{ArBIB} preparation}
\label{subsec:arb-porosity}


If the specimen is mechanically polished (\zcref{sec:mech_pol}), there is a risk of smearing the material into the pores or introducing scratches hindering a pore segmentation.
As detailed in \zcref{sec:bib_pol_moving}, mechanically prepared surfaces of hydrated cements can be improved using \ac{ArBIB} polishing.
However, this still does not allow a close-to-native imaging.
In contrast, \ac{ArBIB} sectioning allows to cut cross-sections (see \zcref{sec:bib_pol_fixed}) of \ac{CSH} needles of hydrated \ac{C3S} (see \zcref{fig:csh-arbib-a} and \zcref{fig:csh-arbib-b}).

\begin{figure}[htb!]
    \centering

	\begin{subfigure}{0.49\textwidth}
		%X:\Mitarbeiter\Florian Kleiner\Einzelbilder\2020_05_19 C3S 28d BIB 6KV 6h\C3S 28d BIB 6KV 6h_006.tif
		\includegraphics[width=\textwidth]{methods/artistic-csh-anim/CSH-cross-section1.pdf}
		\caption{Not embedded specimen. \ac{SE} mode, \qty{2}{\kilo\volt}.}
		\label{fig:csh-arbib-a}
	\end{subfigure}
	\hfill
	\begin{subfigure}{0.49\textwidth}
		%X:\Mitarbeiter\Florian Kleiner\Einzelbilder\selected BIB\Cryo+Harz_055.tif
		\includegraphics[width=\textwidth]{methods/artistic-csh-anim/CSH-cross-section2.pdf}
		\caption{Epoxy resin embedded specimen. \ac{BSE} mode, \qty{1}{\kilo\volt}.}
		\label{fig:csh-arbib-b}
	\end{subfigure}

    \caption{
		\ac{SEM}-images of \ac{ArBIB} cross-sections under cryo conditions of \qty{7}{\day} hydrated \ac{C3S} revealing the cross-section of \ac{CSH}-needles, resembling a star like shape ({\color{blue}blue}).
		Through the lens detector.%
	}
    \label{fig:csh-arbib-SEM}
\end{figure}

These images of a section of a \ac{CSH}-needle illustrate, that the needles are indeed not hollow as suggested by \textcite{double.1976, double.1978} in the 1970s.

As mentioned in \zcref{fig:bib_scheme}, an irradiation volume is expected during the \ac{ArBIB} sectioning.
Since hydrate phases are easily damaged by electron radiation as demonstrated in \zcref{fig:ebeam_damage}, the same could be expected during the ion beam preparation.
To determine the damage induced by this method, a CEM~I was polished at \num{-140} or \qty{20}{\degreeCelsius}.
Comparig images of cryo-prepared (\zcref{fig:csh-arbib-SEM}) areas and areas prepared at room temperature (\zcref{fig:csh-arbib-SEM2}) indicate, that there is no visible damage caused by the ion beam without cooling.

% X:\Mitarbeiter\Florian Kleiner\Einzelbilder\2020_04_29 C3S 28d RT BIB
\begin{figure}[htb!]
    \centering

	\begin{subfigure}{0.49\textwidth}
		\includegraphics[width=\textwidth]{methods/arbib-c3s-rt/fibres.pdf}
		\caption{\ac{CSHo}, open porosity.}
		\label{fig:csh-arbib-c}
	\end{subfigure}
	\hfill
	\begin{subfigure}{0.49\textwidth}
		\includegraphics[width=\textwidth]{methods/arbib-c3s-rt/CHSi_pores.pdf}
		\caption{Close-up of \ac{CSHi} and its porosity.}
		\label{fig:csh-arbib-d}
	\end{subfigure}

    \caption{
		\ac{SEM}-images of cross-sections prepared by \ac{ArBIB} at room temperature of \qty{28}{\day} hydrated \ac{C3S} without epoxy resin intrusion.
		The cross-sections reveal the \ac{CSH}-needle filled capillary pore space (left) and the porosity within \ac{CSHi}.
		Through the lens detector in \ac{SE} mode, \qty{2}{\kilo\volt}.
		}
    \label{fig:csh-arbib-SEM2}
\end{figure}

Similar to the \ac{SEM}-\ac{SE} images of the \textsc{Vycor} samples, the segmentation cannot be done using a global threshold value.
While in \ac{BSE} image a rather distinct grey value range represents a single phase, this is not true for \ac{SE} images.
Therefore, also for these images a local threshold method was applied.
% TODO: results from the ArBIB Paper -> phansalkar thresholding to segment the pores in CSHi

\subsection{\acs*{SEM}}

These images reveal the porosity within \ac{CSHi} and the porosity caused by the interpenetrating network of \ac{CSH}-fibres.
While it is not necessary from a physical point of view, these images also illustrate the advantage of epoxy resin intrusion (\zcref{fig:csh-arbib-b}) into the pore space between the needle structure for image segmentation.
In \zcref{fig:csh-arbib-a, fig:csh-arbib-c} the full pore background and not only a section is visible, making a useful segmentation and evaluation of the pore space unfeasible.
However, the image quality of the nano-scale porosity within \ac{CSHi} is not affected by the epoxy intrusion.
Therefore a segmentation can be done for both preparation variants.

Furthermore, the \ac{SEM}-images show the cross-section of \ac{CSH}-needles, which are star-shaped (\zcref{fig:csh-arbib-b}).
However, this can mainly be observed at the stem of the needle as foils radiate from the centre of the needle, resembling plank roots (buttress roots) of shallowly rooting trees found in tropical forests. % baumhinweis drin lassen ?

This structure can also be found in cracked surfaces of the same material (\zcref{fig:csh-fracture}).
This allows to create a 3D reconstruction of a hydrated \ac{C3S} grain covered in \ac{CSH}-fibres (\zcref{fig:csh-animation}).

\begin{figure}[htb!]
	\centering

	 \begin{subfigure}{0.49\textwidth}
		\includegraphics[width=\textwidth]{methods/artistic-csh-anim/CSH-cross-section.pdf}
	 	\caption{Fractured surface of \qty{28}{\day} hydrated \ac{C3S}.}
	 	\label{fig:csh-fracture}
	 \end{subfigure}
	 \hfill
	 \begin{subfigure}{0.49\textwidth}
		\ifdefined\PRINTABLE
			\includegraphics[width=\textwidth]{methods/artistic-csh-anim/still.png}
		\else
			\animategraphics[palindrome,autoplay,width=\textwidth]{12}{methods/artistic-csh-anim/Test00}{01}{30}
		\fi
		\caption{Artistic 3D reconstruction of \ac{CSHo}.}
		\label{fig:csh-animation}
	 \end{subfigure}

    \caption{
		Fractured surface showing the stems of \ac{CSH} needles and an artistic three-dimensional reconstruction of \ac{CSH} fibres covering a \ac{C3S}-particle based on a \ac{ArBIB} cross-section of hydrated \ac{C3S}. \cite{rossler.2023}
	}
    \label{fig:csh-arbib-ART}
\end{figure}


\FloatBarrier
\subsection{\acs{LTDSC}}

The \ac{LTDSC} measurements of the second freezing pass of \ac{C3S} and \ac{C2S} from \qtyrange{1}{28}{\day} of hydration are shown in \zcref{fig:ltdsc_result_freeze} (upper graphs \ac{C3S}, lower graphs \ac{C2S}).
The graph shows the mean curve from three individual measurements for each specimen age.
The individual measurements can be found in \zcref{fig:ltdsc_result_freeze_raw}.
It was not possible to measure \ac{C2S} after \qty{1}{\day} of hydration, since the material was not hardened enough to be handled and transferred into the \ac{LTDSC} crucible.

% \ifdefined\EXTERNALIZE
% 	\tikzexternaldisable
% \fi
\begin{figure}[!htb]
    \centering

 	\begin{subfigure}{.47\textwidth}
      \begin{tikzpicture}[remember picture]
        \begin{semilogxaxis}[
				blue,
				axis line style={black},
				xmin=1.5, xmax=60,
				ymin=-.20, ymax=0,
				y dir=reverse,
				width=\linewidth,
				height=6cm,
				grid=none,
				xlabel={temperature in \unit{\celsius}},
				xticklabels={\num{-51.2}, \num{-25.3}, \num{-12.4}, \num{-4.6}, \num{-2.0}},
				xtick={2.5,5,10,25,50},
				ytick=\empty,
				scaled y ticks=false,
				scaled x ticks=false,
				axis x line*=top,
            ]

        \end{semilogxaxis}
        \begin{semilogxaxis}[
				legend style={at={(0.95,.7)},anchor=east},
                legend cell align={left},
				xmin=1.5, xmax=60,
				ymin=-.20, ymax=0.07,
				y dir=reverse,
				width=\linewidth,
				height=6cm,
				grid=major,
				xlabel={pore radius $r_\text{P}$ \unit{\nano\meter}},
				ylabel={heat flow $\dot{Q}$ in \unit{\joule\per\second\per\gram}},
				grid style={line width=.1pt, draw=gray!10},
				yticklabel style={
						/pgf/number format/fixed,
						/pgf/number format/precision=1,
						/pgf/number format/fixed zerofill,
				},
				xticklabels={2.5,5,10,25,50},
				xtick={2.5,5,10,25,50},
				scaled y ticks=false,
				scaled x ticks=false,
				clip=false,
            ]
			\addplot[dotted, thick, samples=2, darkgrey,forget plot] coordinates {(3.19,-.2)(3.19,0.07)};
			\node[darkgrey, rotate=90, anchor=west] at (axis cs: 2.51,-.105) {\scriptsize <\,\qty{-40}{\celsius}};

			% std dev of 1d C3S
            \addplot+ [draw=none, no markers, name path=A,forget plot] table [x index=3, col sep=comma,y expr={\thisrowno{1}-\thisrowno{2}}] {figures/ltdsc/reduced/C3S_1d_mean_freeze_segment.csv};
            \addplot+ [draw=none, no markers, name path=B,forget plot] table [x index=3, col sep=comma,y expr={\thisrowno{1}+\thisrowno{2}}] {figures/ltdsc/reduced/C3S_1d_mean_freeze_segment.csv};
            \addplot [fill=yellow, fill opacity=0.2,forget plot] fill between [of=A and B];

			% std dev of 7d C3S
            \addplot+ [draw=none, no markers, name path=A,forget plot] table [x index=3, col sep=comma,y expr={\thisrowno{1}-\thisrowno{2}}] {figures/ltdsc/reduced/C3S_7d_mean_freeze_segment.csv};
            \addplot+ [draw=none, no markers, name path=B,forget plot] table [x index=3, col sep=comma,y expr={\thisrowno{1}+\thisrowno{2}}] {figures/ltdsc/reduced/C3S_7d_mean_freeze_segment.csv};
            \addplot [fill=orange, fill opacity=0.2,forget plot] fill between [of=A and B];

			% std dev of 14d C3S
            \addplot+ [draw=none, no markers, name path=A,forget plot] table [x index=3, col sep=comma,y expr={\thisrowno{1}-\thisrowno{2}}] {figures/ltdsc/reduced/C3S_14d_mean_freeze_segment.csv};
            \addplot+ [draw=none, no markers, name path=B,forget plot] table [x index=3, col sep=comma,y expr={\thisrowno{1}+\thisrowno{2}}] {figures/ltdsc/reduced/C3S_14d_mean_freeze_segment.csv};
            \addplot [fill=red, fill opacity=0.2,forget plot] fill between [of=A and B];

			% std dev of 21d C3S
            \addplot+ [draw=none, no markers, name path=A,forget plot] table [x index=3, col sep=comma,y expr={\thisrowno{1}-\thisrowno{2}}] {figures/ltdsc/reduced/C3S_21d_mean_freeze_segment.csv};
            \addplot+ [draw=none, no markers, name path=B,forget plot] table [x index=3, col sep=comma,y expr={\thisrowno{1}+\thisrowno{2}}] {figures/ltdsc/reduced/C3S_21d_mean_freeze_segment.csv};
            \addplot [fill=green, fill opacity=0.2,forget plot] fill between [of=A and B];

			% std dev of 28d C3S
            \addplot+ [draw=none, no markers, name path=A,forget plot] table [x index=3, col sep=comma,y expr={\thisrowno{1}-\thisrowno{2}}] {figures/ltdsc/reduced/C3S_28d_mean_freeze_segment.csv};
            \addplot+ [draw=none, no markers, name path=B,forget plot] table [x index=3, col sep=comma,y expr={\thisrowno{1}+\thisrowno{2}}] {figures/ltdsc/reduced/C3S_28d_mean_freeze_segment.csv};
            \addplot [fill=blue, fill opacity=0.2,forget plot] fill between [of=A and B];


            \addplot [ no marks, yellow] table [x index=3, y index=1, col sep=comma] {figures/ltdsc/reduced/C3S_1d_mean_freeze_segment.csv};
            %\addlegendentry{\ac{C3S}, \qty{1}{\day}};
            \addplot [ no marks, orange] table [x index=3, y index=1, col sep=comma] {figures/ltdsc/reduced/C3S_7d_mean_freeze_segment.csv};
            %\addlegendentry{\ac{C3S}, \qty{7}{\day}};
            \addplot [ no marks, red] table [x index=3, y index=1, col sep=comma] {figures/ltdsc/reduced/C3S_14d_mean_freeze_segment.csv};
            %\addlegendentry{\ac{C3S}, \qty{14}{\day}};
            \addplot [ no marks, green] table [x index=3, y index=1, col sep=comma] {figures/ltdsc/reduced/C3S_21d_mean_freeze_segment.csv};
            %\addlegendentry{\ac{C3S}, \qty{21}{\day}};
            \addplot [ no marks, blue] table [x index=3, y index=1, col sep=comma] {figures/ltdsc/reduced/C3S_28d_mean_freeze_segment.csv};
            %\addlegendentry{\ac{C3S}, \qty{28}{\day}};

			\draw[draw=black, densely dotted] (2.75,-.11) rectangle (1.75,-0.03);
			\node[draw=none] (top-3) at (axis cs:2.75,-.11) {};
			\node[draw=none] (bottom-3) at (axis cs:2.75,-.03) {};

            \node[anchor=west, align=left] at (.8,-0.25) {\textbf{\ac{C3S}}};
        \end{semilogxaxis}
	  \end{tikzpicture}

	\end{subfigure}
	\hfill
	\begin{subfigure}{.47\textwidth}

	  \begin{tikzpicture}[remember picture]
        \begin{axis}[
			blue,
			axis line style={black},
			xmin=1.75, xmax=2.75,
			ymin=-.11, ymax=-0.03,
			y dir=reverse,
			width=.9\linewidth,
			height=6cm,
			grid=none,
			xlabel={temperature in \unit{\celsius}},
			xticklabels={\num{-73.3}, \num{-64.1}, \num{-56.9}, \num{-51.2}, \num{-46.5}},
			xtick={1.75,2.00,2.25,2.50,2.75},
			ytick=\empty,
			scaled y ticks=false,
			scaled x ticks=false,,
			axis x line*=top
		]

		\end{axis}
		\begin{axis}[
				legend style={at={(0.97,.97)},anchor=north east},
				legend cell align={left},
				xmin=1.75, xmax=2.75,
				ymin=-.11, ymax=-0.03,
				y dir=reverse,
				width=.9\linewidth,
				height=6cm,
				grid=major,
				xlabel={pore radius $r_\text{P}$ \unit{\nano\meter}},
				ylabel={heat flow $\dot{Q}$ in \unit{\joule\per\second\per\gram}},
				grid style={line width=.1pt, draw=gray!10},
				xticklabels={1.75,2.00,2.25,2.50,2.75},
				xtick={1.75,2.00,2.25,2.50,2.75},
				yticklabel style={
						/pgf/number format/fixed,
						/pgf/number format/precision=2,
						/pgf/number format/fixed zerofill,
				},
				scaled y ticks=false,
				scaled x ticks=false,
				yticklabel pos=right,
			]

			% std dev of 1d C3S
            \addplot+ [draw=none, no markers, name path=A,forget plot] table [x index=3, col sep=comma,y expr={\thisrowno{1}-\thisrowno{2}}] {figures/ltdsc/reduced/zoom/C3S_1d_mean_freeze_segment.csv};
            \addplot+ [draw=none, no markers, name path=B,forget plot] table [x index=3, col sep=comma,y expr={\thisrowno{1}+\thisrowno{2}}] {figures/ltdsc/reduced/zoom/C3S_1d_mean_freeze_segment.csv};
            \addplot [fill=yellow, fill opacity=0.2,forget plot] fill between [of=A and B];

			% std dev of 7d C3S
            \addplot+ [draw=none, no markers, name path=A,forget plot] table [x index=3, col sep=comma,y expr={\thisrowno{1}-\thisrowno{2}}] {figures/ltdsc/reduced/zoom/C3S_7d_mean_freeze_segment.csv};
            \addplot+ [draw=none, no markers, name path=B,forget plot] table [x index=3, col sep=comma,y expr={\thisrowno{1}+\thisrowno{2}}] {figures/ltdsc/reduced/zoom/C3S_7d_mean_freeze_segment.csv};
            \addplot [fill=orange, fill opacity=0.2,forget plot] fill between [of=A and B];

			% std dev of 14d C3S
            \addplot+ [draw=none, no markers, name path=A,forget plot] table [x index=3, col sep=comma,y expr={\thisrowno{1}-\thisrowno{2}}] {figures/ltdsc/reduced/zoom/C3S_14d_mean_freeze_segment.csv};
            \addplot+ [draw=none, no markers, name path=B,forget plot] table [x index=3, col sep=comma,y expr={\thisrowno{1}+\thisrowno{2}}] {figures/ltdsc/reduced/zoom/C3S_14d_mean_freeze_segment.csv};
            \addplot [fill=red, fill opacity=0.2,forget plot] fill between [of=A and B];

			% std dev of 21d C3S
            \addplot+ [draw=none, no markers, name path=A,forget plot] table [x index=3, col sep=comma,y expr={\thisrowno{1}-\thisrowno{2}}] {figures/ltdsc/reduced/zoom/C3S_21d_mean_freeze_segment.csv};
            \addplot+ [draw=none, no markers, name path=B,forget plot] table [x index=3, col sep=comma,y expr={\thisrowno{1}+\thisrowno{2}}] {figures/ltdsc/reduced/zoom/C3S_21d_mean_freeze_segment.csv};
            \addplot [fill=green, fill opacity=0.2,forget plot] fill between [of=A and B];

			% std dev of 28d C3S
            \addplot+ [draw=none, no markers, name path=A,forget plot] table [x index=3, col sep=comma,y expr={\thisrowno{1}-\thisrowno{2}}] {figures/ltdsc/reduced/zoom/C3S_28d_mean_freeze_segment.csv};
            \addplot+ [draw=none, no markers, name path=B,forget plot] table [x index=3, col sep=comma,y expr={\thisrowno{1}+\thisrowno{2}}] {figures/ltdsc/reduced/zoom/C3S_28d_mean_freeze_segment.csv};
            \addplot [fill=blue, fill opacity=0.2,forget plot] fill between [of=A and B];

			% mean values of the 1d C3S
			\addplot+[no marks, orange] table [x index=3, y index=1, col sep=comma] {figures/ltdsc/reduced/zoom/C3S_1d_mean_freeze_segment.csv};
			\addlegendentry{\qty{1}{\day}};
			% mean values of the 7d C3S
			\addplot+[no marks, orange] table [x index=3, y index=1, col sep=comma] {figures/ltdsc/reduced/zoom/C3S_7d_mean_freeze_segment.csv};
			\addlegendentry{\qty{7}{\day}};
			% mean values of the 14d C3S
			\addplot+[no marks, red] table [x index=3, y index=1, col sep=comma] {figures/ltdsc/reduced/zoom/C3S_14d_mean_freeze_segment.csv};
			\addlegendentry{\qty{14}{\day}};
			% mean values of the 21d C3S
			\addplot+[no marks, green] table [x index=3, y index=1, col sep=comma] {figures/ltdsc/reduced/zoom/C3S_21d_mean_freeze_segment.csv};
			\addlegendentry{\qty{21}{\day}};
			% mean values of the 28d C3S
			\addplot+[no marks, blue] table [x index=3, y index=1, col sep=comma] {figures/ltdsc/reduced/zoom/C3S_28d_mean_freeze_segment.csv};
			\addlegendentry{\qty{28}{\day}};

            %\node[anchor=west, align=left] at (axis cs:2,-0.09) {\ac{C3S}};

			\node[draw=none] (top-4) at (axis cs:1.75,-.11) {};
			\node[draw=none] (bottom-4) at (axis cs:1.75,-.03) {};
		\end{axis}
	  \end{tikzpicture}
	\end{subfigure}
    \begin{tikzpicture}[remember picture,overlay]
        \draw[dashed,grey] (top-3) -- (top-4);
        \draw[dashed,grey] (bottom-3) -- (bottom-4);
    \end{tikzpicture}

	\vspace{1\baselineskip}

 	\begin{subfigure}{.47\textwidth}
      \begin{tikzpicture}[remember picture]
        \begin{semilogxaxis}[
				blue,
				axis line style={black},
				xmin=1.5, xmax=60,
				ymin=-.20, ymax=0,
				y dir=reverse,
				width=\linewidth,
				height=6cm,
				grid=none,
				xlabel={temperature in \unit{\celsius}},
				xticklabels={\num{-51.2}, \num{-25.3}, \num{-12.4}, \num{-4.6}, \num{-2.0}},
				xtick={2.5,5,10,25,50},
				ytick=\empty,
				scaled y ticks=false,
				scaled x ticks=false,,
				axis x line*=top
            ]

        \end{semilogxaxis}
        \begin{semilogxaxis}[
				legend style={at={(0.95,.7)},anchor=east},
                legend cell align={left},
				xmin=1.5, xmax=60,
				ymin=-.20, ymax=0.07,
				y dir=reverse,
				width=\linewidth,
				height=6cm,
				grid=major,
				xlabel={pore radius $r_\text{P}$ \unit{\nano\meter}},
				ylabel={heat flow $\dot{Q}$ in \unit{\joule\per\second\per\gram}},
				grid style={line width=.1pt, draw=gray!10},
				yticklabel style={
						/pgf/number format/fixed,
						/pgf/number format/fixed zerofill,
						/pgf/number format/precision=1
				},
				xticklabels={2.5,5,10,25,50},
				xtick={2.5,5,10,25,50},
				scaled y ticks=false,
				scaled x ticks=false,
				clip=false,
            ]
			\addplot[dotted, thick, samples=2, darkgrey,forget plot] coordinates {(3.19,-.2)(3.19,0.07)};
			\node[darkgrey, rotate=90, anchor=west] at (axis cs: 2.51,-.105) {\scriptsize <\,\qty{-40}{\celsius}};

			% std dev of 7d C2S
            \addplot+ [draw=none, no markers, name path=A,forget plot] table [x index=3, col sep=comma,y expr={\thisrowno{1}-\thisrowno{2}}] {figures/ltdsc/reduced/C2S_7d_mean_freeze_segment.csv};
            \addplot+ [draw=none, no markers, name path=B,forget plot] table [x index=3, col sep=comma,y expr={\thisrowno{1}+\thisrowno{2}}] {figures/ltdsc/reduced/C2S_7d_mean_freeze_segment.csv};
            \addplot [fill=orange, fill opacity=0.2,forget plot] fill between [of=A and B];

			% std dev of 14d C2S
            \addplot+ [draw=none, no markers, name path=A,forget plot] table [x index=3, col sep=comma,y expr={\thisrowno{1}-\thisrowno{2}}] {figures/ltdsc/reduced/C2S_14d_mean_freeze_segment.csv};
            \addplot+ [draw=none, no markers, name path=B,forget plot] table [x index=3, col sep=comma,y expr={\thisrowno{1}+\thisrowno{2}}] {figures/ltdsc/reduced/C2S_14d_mean_freeze_segment.csv};
            \addplot [fill=red, fill opacity=0.2,forget plot] fill between [of=A and B];

			% std dev of 21d C2S
            \addplot+ [draw=none, no markers, name path=A,forget plot] table [x index=3, col sep=comma,y expr={\thisrowno{1}-\thisrowno{2}}] {figures/ltdsc/reduced/C2S_21d_mean_freeze_segment.csv};
            \addplot+ [draw=none, no markers, name path=B,forget plot] table [x index=3, col sep=comma,y expr={\thisrowno{1}+\thisrowno{2}}] {figures/ltdsc/reduced/C2S_21d_mean_freeze_segment.csv};
            \addplot [fill=green, fill opacity=0.2,forget plot] fill between [of=A and B];

			% std dev of 28d C2S
            \addplot+ [draw=none, no markers, name path=A,forget plot] table [x index=3, col sep=comma,y expr={\thisrowno{1}-\thisrowno{2}}] {figures/ltdsc/reduced/C2S_28d_mean_freeze_segment.csv};
            \addplot+ [draw=none, no markers, name path=B,forget plot] table [x index=3, col sep=comma,y expr={\thisrowno{1}+\thisrowno{2}}] {figures/ltdsc/reduced/C2S_28d_mean_freeze_segment.csv};
            \addplot [fill=blue, fill opacity=0.2,forget plot] fill between [of=A and B];


            \addplot [ no marks, orange] table [x index=3, y index=1, col sep=comma] {figures/ltdsc/reduced/C2S_7d_mean_freeze_segment.csv};
            %\addlegendentry{\ac{C2S}, \qty{7}{\day}};
            \addplot [ no marks, red] table [x index=3, y index=1, col sep=comma] {figures/ltdsc/reduced/C2S_14d_mean_freeze_segment.csv};
            %\addlegendentry{\ac{C2S}, \qty{14}{\day}};
            \addplot [ no marks, green] table [x index=3, y index=1, col sep=comma] {figures/ltdsc/reduced/C2S_21d_mean_freeze_segment.csv};
            %\addlegendentry{\ac{C2S}, \qty{21}{\day}};
            \addplot [ no marks, blue] table [x index=3, y index=1, col sep=comma] {figures/ltdsc/reduced/C2S_28d_mean_freeze_segment.csv};
            %\addlegendentry{\ac{C2S}, \qty{28}{\day}};

			\draw[draw=black, densely dotted] (2.75,-.11) rectangle (1.75,-0.03);
			\node[draw=none] (top-1) at (axis cs:2.75,-.11) {};
			\node[draw=none] (bottom-1) at (axis cs:2.75,-.03) {};

            \node[anchor=west, align=left] at (.8,-0.25) {\textbf{\ac{C2S}}};

        \end{semilogxaxis}
	  \end{tikzpicture}

	\end{subfigure}
	\hfill
	\begin{subfigure}{.47\textwidth}

	  \begin{tikzpicture}[remember picture]
        \begin{axis}[
			blue,
			axis line style={black},
			xmin=1.75, xmax=2.75,
			ymin=-.11, ymax=-0.03,
			y dir=reverse,
			width=.9\linewidth,
			height=6cm,
			grid=none,
			xlabel={temperature in \unit{\celsius}},
			xticklabels={\num{-73.3}, \num{-64.1}, \num{-56.9}, \num{-51.2}, \num{-46.5}},
			xtick={1.75,2.00,2.25,2.50,2.75},
			ytick=\empty,
			scaled y ticks=false,
			scaled x ticks=false,,
			axis x line*=top
		]

		\end{axis}
		\begin{axis}[
				legend style={at={(0.97,.97)},anchor=north east},
				legend cell align={left},
				xmin=1.75, xmax=2.75,
				ymin=-.11, ymax=-0.03,
				y dir=reverse,
				width=.9\linewidth,
				height=6cm,
				grid=major,
				xlabel={pore radius $r_\text{P}$ \unit{\nano\meter}},
				ylabel={heat flow $\dot{Q}$ in \unit{\joule\per\second\per\gram}},
				grid style={line width=.1pt, draw=gray!10},
				xticklabels={1.75,2.00,2.25,2.50,2.75},
				xtick={1.75,2.00,2.25,2.50,2.75},
				yticklabel style={
						/pgf/number format/fixed,
						/pgf/number format/precision=2,
						/pgf/number format/fixed zerofill,
				},
				scaled y ticks=false,
				scaled x ticks=false,
				yticklabel pos=right,
			]
			% std dev of 7d C2S
            \addplot+ [draw=none, no markers, name path=A,forget plot] table [x index=3, col sep=comma,y expr={\thisrowno{1}-\thisrowno{2}}] {figures/ltdsc/reduced/zoom/C2S_7d_mean_freeze_segment.csv};
            \addplot+ [draw=none, no markers, name path=B,forget plot] table [x index=3, col sep=comma,y expr={\thisrowno{1}+\thisrowno{2}}] {figures/ltdsc/reduced/zoom/C2S_7d_mean_freeze_segment.csv};
            \addplot [fill=orange, fill opacity=0.2,forget plot] fill between [of=A and B];

			% std dev of 14d C2S
            \addplot+ [draw=none, no markers, name path=A,forget plot] table [x index=3, col sep=comma,y expr={\thisrowno{1}-\thisrowno{2}}] {figures/ltdsc/reduced/zoom/C2S_14d_mean_freeze_segment.csv};
            \addplot+ [draw=none, no markers, name path=B,forget plot] table [x index=3, col sep=comma,y expr={\thisrowno{1}+\thisrowno{2}}] {figures/ltdsc/reduced/zoom/C2S_14d_mean_freeze_segment.csv};
            \addplot [fill=red, fill opacity=0.2,forget plot] fill between [of=A and B];

			% std dev of 21d C2S
            \addplot+ [draw=none, no markers, name path=A,forget plot] table [x index=3, col sep=comma,y expr={\thisrowno{1}-\thisrowno{2}}] {figures/ltdsc/reduced/zoom/C2S_21d_mean_freeze_segment.csv};
            \addplot+ [draw=none, no markers, name path=B,forget plot] table [x index=3, col sep=comma,y expr={\thisrowno{1}+\thisrowno{2}}] {figures/ltdsc/reduced/zoom/C2S_21d_mean_freeze_segment.csv};
            \addplot [fill=green, fill opacity=0.2,forget plot] fill between [of=A and B];

			% std dev of 28d C2S
            \addplot+ [draw=none, no markers, name path=A,forget plot] table [x index=3, col sep=comma,y expr={\thisrowno{1}-\thisrowno{2}}] {figures/ltdsc/reduced/zoom/C2S_28d_mean_freeze_segment.csv};
            \addplot+ [draw=none, no markers, name path=B,forget plot] table [x index=3, col sep=comma,y expr={\thisrowno{1}+\thisrowno{2}}] {figures/ltdsc/reduced/zoom/C2S_28d_mean_freeze_segment.csv};
            \addplot [fill=blue, fill opacity=0.2,forget plot] fill between [of=A and B];

			% mean values of the 7d C2S
			\addplot+[no marks, orange] table [each nth point=3, filter discard warning=false, unbounded coords=discard, x index=3, y index=1, col sep=comma] {figures/ltdsc/reduced/zoom/C2S_7d_mean_freeze_segment.csv};
			\addlegendentry{\qty{7}{\day}};
			% mean values of the 14d C2S
			\addplot+[no marks, red] table [each nth point=3, filter discard warning=false, unbounded coords=discard,x index=3, y index=1, col sep=comma] {figures/ltdsc/reduced/zoom/C2S_14d_mean_freeze_segment.csv};
			\addlegendentry{\qty{14}{\day}};
			% mean values of the 21d C2S
			\addplot+[no marks, green] table [each nth point=3, filter discard warning=false, unbounded coords=discard,x index=3, y index=1, col sep=comma] {figures/ltdsc/reduced/zoom/C2S_21d_mean_freeze_segment.csv};
			\addlegendentry{\qty{21}{\day}};
			% mean values of the 28d C2S
			\addplot+[no marks, blue] table [each nth point=3, filter discard warning=false, unbounded coords=discard,x index=3, y index=1, col sep=comma] {figures/ltdsc/reduced/zoom/C2S_28d_mean_freeze_segment.csv};
			\addlegendentry{\qty{28}{\day}};


			\node[draw=none] (top-2) at (axis cs:1.75,-.11) {};
			\node[draw=none] (bottom-2) at (axis cs:1.75,-.03) {};
		\end{axis}
	  \end{tikzpicture}
	\end{subfigure}
    \begin{tikzpicture}[remember picture,overlay]
        \draw[dashed,grey] (top-1) -- (top-2);
        \draw[dashed,grey] (bottom-1) -- (bottom-2);
    \end{tikzpicture}

    \caption{
		\ac{LTDSC} measurements (mean of three measurements) for the second freezing pass of \ac{C3S} (top) and \ac{C2S} (bottom) from \qtyrange{1}{28}{\day} of hydration.
		Every specimen age was measured using three specimens.
		On the left-hand side, the full measurement is plotted (pore radius and temperature is plotted logarithmic).
		On the right-hand side is a magnified section (linear plot), highlighting the data correlating to the gel pores.
		Note that pore sizes below \qty{3.19}{\nano\meter} (<\,\qty{-40}{\celsius}) are undefined as noted in \zcref{eq:brun_freeze}.
	}
    \label{fig:ltdsc_result_freeze}
\end{figure}

The graph on the right is a magnified section of the left graph, showing the peaks correlating to the gel pores.
Applying \zcref{eq:brun_freeze}, a pore radius of \qtyrange{2.00}{2.25}{\nano\meter} can be calculated.
However, according to \textcite{brun.1977} this equation is not valid for temperatures below \qty{-40}{\celsius}.
The figure shows that the heat flow signal of the gel pores is increasing with the age of the specimen, which indicates an increase in the gel pore volume with increasing specimen age.

On first glance, the \ac{C3S} results look similar to the \ac{C2S} results.
However, they have a lower signal response in the temperature range corresponding to the gel pores.
This indicates a lower gel pore volume for \ac{C3S} compared to \ac{C2S}.
Furthermore, the peaks are slightly shifted to smaller pore radii indicating smaller gel pores.

However, the freezing curve cannot be used to determine the volume of the gel pores.
In order to calculate the ice mass within the gel pores and thus also their volume, the heat flow signal of the two melting curves of every experiment was integrated over the measured temperature ranges and subtracted from each other as shown in \zcref{eq:ltdsc_ice_integral, eq:ltdsc_ice2}.
In order to do that, a base line correction was applied to the raw data.
\zcref{fig:ltdsc_result_c2s_melt} illustrates this process for the results of \qty{28}{\day} hydrated \ac{C2S}.
%{\color{red}The detailed procedure is published on github \cite{kleiner.2025} %\url{https://github.com/kleinerELM/Kalorimeterauswertung}
%}

The raw data are plotted as dashed lines, while the baseline corrected data are plotted as solid lines.
The baseline functions $\dot{Q}_\text{B, P1}(T)$ and $\dot{Q}_\text{B, P2}(T)$ are plotted in the graph.
The areas under the solid curves are marked in red and yellow.
The red area represents the gel pore volume, while the yellow area represents the capillary pore volume.

\begin{figure}[!htb]
    \centering

    \begin{tikzpicture}
        \begin{axis}[
				legend style={at={(0.1,.95)},anchor=north west},
                legend cell align={left},
				xmin=-60, xmax=15,
				ymin=0, ymax=0.6,
				width=\linewidth,
				height=7cm,
				grid=major,
				xlabel={temperature $\vartheta$ in \unit{\celsius}},
				ylabel={heat flow $\dot{Q}$ in \unit{\joule\per\second\per\gram}},
				xtick={-60, -50,...,10},
				grid style={line width=.1pt, draw=gray!10},
				yticklabel style={
						/pgf/number format/fixed,
						/pgf/number format/fixed zerofill,
						/pgf/number format/precision=1
				},
				scaled y ticks=false,
				scaled x ticks=false
            ]


            \addplot[fill=red, fill opacity=0.25, draw=none] coordinates {
                (-57, 0) (-56,0) (-56, 0.6) (-57,0.6) (-57,0)
            };
            \addplot[fill=red, fill opacity=0.25, draw=none] coordinates {
                (-21.5, 0) (-20.5,0) (-20.5, 0.6) (-21.5,0.6) (-21.5,0)
            };

            \addplot[fill=red, fill opacity=0.25, draw=none] coordinates {
                (9, 0) (10,0) (10, 0.6) (9,0.6) (9,0)
            };

            \node[anchor=west, align=left, font=\footnotesize]
                at (axis cs:-53,0.225) {
					$\dot{Q}_\text{B, P1}(T) = 0.00091T + 0.04939$\\
					$\dot{Q}_\text{B, P2}(T) = 0.00073T + 0.04064$};

            \node[anchor=west, align=left] at (axis cs:-4,0.5) {A};
            \node[anchor=west, align=left] at (axis cs:3.5,0.55) {B};

            \addplot [ no marks, green, dashed] table [each nth point=4, filter discard warning=false, unbounded coords=discard,x index=0, y index=1, col sep=comma] {figures/ltdsc/ExpDat_C2S_28d-1_melt_segment_8.csv};
            \addlegendentry{P1, raw data};

            \addplot [ no marks, blue, dashed] table [each nth point=4, filter discard warning=false, unbounded coords=discard,x index=0, y index=1, col sep=comma] {figures/ltdsc/ExpDat_C2S_28d-1_melt_segment_16.csv};
            \addlegendentry{P2, raw data};


			\addplot[thick, domain=0:1, color=blue, fill = yellow, fill opacity=0.3, forget plot] table [each nth point=4, filter discard warning=false, unbounded coords=discard,x index=0, y index=2, col sep=comma] {figures/ltdsc/ExpDat_C2S_28d-1_melt_segment_16.csv};

			\addplot[domain=0:1, color=green, fill = red, fill opacity=0.3, forget plot] table [each nth point=4, filter discard warning=false, unbounded coords=discard,x index=0, y index=2, col sep=comma] {figures/ltdsc/ExpDat_C2S_28d-1_melt_segment_8.csv};


            \addplot [thick, no marks, green] table [each nth point=4, filter discard warning=false, unbounded coords=discard,x index=0, y index=2, col sep=comma] {figures/ltdsc/ExpDat_C2S_28d-1_melt_segment_8.csv};
            \addlegendentry{P1, baseline corrected};

            \addplot [thick, no marks, blue] table [each nth point=4, filter discard warning=false, unbounded coords=discard,x index=0, y index=2, col sep=comma] {figures/ltdsc/ExpDat_C2S_28d-1_melt_segment_16.csv};
            \addlegendentry{P2, baseline corrected};

        \end{axis}
	\end{tikzpicture}

    \caption{
		\ac{LTDSC} curves for the first and second melting pass of \qty{28}{\day} hydrated \ac{C2S}.
		The red area correlate to the capillary porosity, while the yellow area correlates to the gel porosity.
		The value range used for the linear baseline regression is marked red between \qtyrange{-57}{-56}{\celsius} (P2) or \qtyrange{-21.5}{-20.5}{\celsius} (P1) and \qtyrange{9}{10}{\celsius} (P1 and P2) determine .
	}
    \label{fig:ltdsc_result_c2s_melt}
\end{figure}



The melting curve shows two distinctive peaks, one at \qty{0}{\celsius} (A) and one at \qty{2}{\celsius} (B).
This indicates, that the heating rate of \qty{2}{\kelvin\per\minute} was higher than the inertia of the melting process, as the ice should melt at or before \qty{0}{\celsius}, but is delayed due to the fast temperature change (\qty{2}{\kelvin\per\minute}).

Therefore, an experiment comparing with a slower heating rate (\qty{0.5}{\kelvin\per\minute}) was carried out (\zcref{fig:ltdsc_result_heating_rate_a}).
This results in a shift of the melting peak in the direction of \qty{0}{\celsius} and a seemingly smaller heat output.
However, difference in peak height is caused by a change in the measuring interval.
To mitigate this, the dashed curve indicates the heat flow $\dot{Q}$ of the \qty{0.5}{\kelvin\per\minute} heat rate, corrected for the interval difference . %TODO: ist das legitim?

\begin{figure}[!htb]
    \centering
	\begin{subfigure}{.485\textwidth}
      \begin{tikzpicture}
        \begin{axis}[
			legend style={at={(0.05,.95)},anchor=north west,font=\footnotesize},
                legend cell align={left},
				xmin=-50, xmax=10,
				ymin=-0.01, ymax=0.2,
				width=\linewidth,
				height=7cm,
				grid=major,
				xlabel={temperature $\vartheta$ in \unit{\celsius}},
				ylabel={heat flow $\dot{Q}$ in \unit{\joule\per\second\per\gram}},
				grid style={line width=.1pt, draw=gray!10},
				yticklabel style={
						/pgf/number format/fixed,
						/pgf/number format/fixed zerofill,
						/pgf/number format/precision=2
				},
				scaled y ticks=false,
				scaled x ticks=false
            ]

			\addplot[samples=2, darkgrey,forget plot] coordinates {(15,0)(-60,0)};
			\addplot[samples=2, darkgrey,forget plot] coordinates {(0,-.01)(0,.2)};

            \addplot [thick, no marks, blue] table [each nth point=4, filter discard warning=false, unbounded coords=discard,x index=0, y index=2, col sep=comma] {figures/ltdsc/ExpDat_M28n2-28d_melt_segment_9.csv};
            \addlegendentry{\qty{0.5}{\kelvin\per\minute}};

			\addplot [dashed, no marks, blue] table [each nth point=4, filter discard warning=false, unbounded coords=discard, x index=0, y expr=\thisrowno{2}*4, col sep=comma] {figures/ltdsc/ExpDat_M28n2-28d_melt_segment_9.csv};
			\addlegendentry{\qty{0.5}{\kelvin\per\minute}$\,\times\,4$};


            \addplot [thick, no marks, red] table [each nth point=4, filter discard warning=false, unbounded coords=discard,x index=0, y index=2, col sep=comma] {figures/ltdsc/ExpDat_M28n2-28d_melt_segment_18.csv};
            \addlegendentry{\qty{2.0}{\kelvin\per\minute}};
        \end{axis}
	  \end{tikzpicture}
	  \caption{Melting heat flow}
	  \label{fig:ltdsc_result_heating_rate_a}

	\end{subfigure}
	\hfill
	\begin{subfigure}{.485\textwidth}

	  \begin{tikzpicture}[]
		\begin{axis}[
				legend style={at={(0.05,.95)},anchor=north west,font=\footnotesize},
				legend cell align={left},
				xmin=-50, xmax=10,
				ymin=-10, ymax=180,
				width=\linewidth,
				height=7cm,
				grid=major,
				xlabel={temperature $\vartheta$ in \unit{\celsius}},
				ylabel={cum. ice volume in \unit{\micro\liter\per\gram}},
				grid style={line width=.1pt, draw=gray!10},
				yticklabel style={
					/pgf/number format/fixed,
					/pgf/number format/precision=0
				},
				scaled y ticks=false,
				scaled x ticks=false,
				yticklabel pos=right,
			]
			\addplot[samples=2, darkgrey,forget plot] coordinates {(15,0)(-60,0)};

            \addplot [no marks, cyan] table [each nth point=10, filter discard warning=false, unbounded coords=discard,x index=0, y expr={\thisrowno{4}*1000}, col sep=comma] {figures/ltdsc/ExpDat_M28n2-28d_ice_9.csv};
            \addlegendentry{\qty{0.5}{\kelvin\per\minute}};

            \addplot [no marks, orange] table [each nth point=10, filter discard warning=false, unbounded coords=discard,x index=0,y expr={\thisrowno{4}*1000}, col sep=comma] {figures/ltdsc/ExpDat_M28n2-28d_ice_18.csv};
            \addlegendentry{\qty{2.0}{\kelvin\per\minute}};

		\end{axis}
	  \end{tikzpicture}
	  \caption{Cumulative ice volume}
	  \label{fig:ltdsc_result_heating_rate_b}

	\end{subfigure}

    \caption{\ac{LTDSC} curves and the cumulative ice volume per specimen weight for a slow (\qty{0.5}{\kelvin\per\minute}) and fast (\qty{2}{\kelvin\per\minute}) heating curves from \qtyrange{-50}{15}{\celsius} of \qty{28}{\day} hydrated cement (\ac{wb} = \num{0.5}).}
    \label{fig:ltdsc_result_heating_rate}
\end{figure}

The ice mass should not be affected by the heating rate.
The simplified approach to calculate the ice volume as used by \textcite{mueller.2021,unbehau.2025} (\zcref{eq:ltdsc_ice_integral}) gives an ice volume of \qty{168.2}{\micro\liter\per\gram} (\qty{0.5}{\kelvin\per\minute}) and \qty{154.8}{\micro\liter\per\gram} (\qty{2}{\kelvin\per\minute}), respectively.
However, the calculation of the ice mass is heavily influenced by the way, this baseline of the heat flow $Q_\text{B}$ is determined.
Since the temperature ranges to process the baseline correction are selected manually, results from this approach may differ depending on the person processing the data.
\zcref{fig:basle_line_ice_vol} illustrates the influence on the calculated ice volume, which can differ $\pm$\qty{10}{\micro\liter\per\gram} due to different selected temperature ranges to obtain the baseline function.
Therefore, the temperature ranges used for this correction were held constant for all further calculations.
This keeps different measurements, run with the same heating rate, comparable.

% The cumulative ice mass calculated using both methods is shown in \zcref{fig:ltdsc_result_heating_rate_b}.
% As can be seen in the zoomed in section, the manual approach gives more plausible curves, starting at roughly \qty{-47}{\celsius}, while the automatic approach seems to loose ice mass until \num{-38} ore \qty{-34}{\celsius}.
% This is caused by the intersection point of the base line corrected raw data with the x-axis.

While a slower heating rate causes improved results, the time required for a single measurement was deemed unpractical.
%It would not have been possible to obtain three measurements for every hydration time without a significant change in sample age between repeats.
Since the specimens are still wet and still hydrating, the repeated measurements had to be carried out in a timely manner.
Therefore, it has been decided to use a faster heating rate of \qty{2}{\kelvin\per\minute}. % and the automatic approach, which should cause a comparable error within the same measurement regime.

\begin{figure}[htb]
    \centering
	%\tikzsetnextfilename{ltdsc_result_C2S_melt_bar} % Gibt den Dateinamen der externen Grafik vor
	\begin{subfigure}[t]{.49\textwidth}
	  \begin{tikzpicture}
		\begin{axis}[
			width  = \textwidth,
			height = 8cm,
			major x tick style = transparent,
			ybar=2*\pgflinewidth,
			ymin=0, ymax=370,
			bar width=12pt,
			ymajorgrids = true,
			ylabel = {cum. ice volume in \unit{\micro\liter\per\gram}},
			xticklabels={
				\qty{7}{\day},
				\qty{14}{\day},
				\qty{21}{\day},
				\qty{28}{\day}
			},
			xtick = data,
			scaled y ticks = false,
			enlarge x limits=0.15,
			ymin=0,
			legend cell align=left,
			legend style={
					at={(0.98,0.97)},
					anchor=north east,
					column sep=4pt
			},
			legend columns=-1,
			scaled y ticks = false,
			error bars/y dir=both,
			error bars/y explicit,
			ytick={0,100,200,300,400},
			visualization depends on={\thisrow{pores_b_std} \as \offset},
			nodes near coords,
			node near coords style={
				shift={(axis direction cs:0,\offset)},
				rotate=90,
				anchor=west,
				/pgf/number format/.cd,fixed zerofill,precision=1,
			},
		]
			\addplot[
				orange,fill=orange,mark=none,
				error bars/.cd,
				error bar style={black},
			] table [ x expr=\coordindex, y=pores_a_mean, y error=pores_a_std, col sep=comma] {figures/ltdsc/C2S_ice_vol_dev.csv};

			\addplot[
				cyan,fill=cyan,mark=none,
				error bars/.cd,
				error bar style={black}
			] table [ x expr=\coordindex, y=pores_b_mean, y error=pores_b_std, col sep=comma] {figures/ltdsc/C2S_ice_vol_dev.csv};

			% \addplot[
			% 	blue,fill=blue,mark=none,
			% 	error bars/.cd,
			% 	error bar style={black},
			% ] table [ x expr=\coordindex, y=pores_c_mean, y error=pores_c_std, col sep=comma] {figures/ltdsc/C2S_ice_vol_dev.csv};

			% \legend{gel pores, capillary pores, total pores}
		\end{axis}
	  \end{tikzpicture}

      \caption{Ice volume of \ac{C2S}}
 	  \label{fig:ltdsc_result_C2S_melt_bar}
	\end{subfigure}
	\hfill
	\begin{subfigure}[t]{.49\textwidth}
	  \begin{tikzpicture}
		\begin{axis}[
			width  = \textwidth,
			height = 8cm,
			major x tick style = transparent,
			ybar=2*\pgflinewidth,
			ymin=0, ymax=370,
			bar width=12pt,
			ymajorgrids = true,
			ylabel = {ice volume in \unit{\micro\liter\per\gram}},
			xticklabels={
				\qty{1}{\day},
				\qty{7}{\day},
				\qty{14}{\day},
				\qty{21}{\day},
				\qty{28}{\day}
			},
			xtick = data,
			enlarge x limits=0.15,
			ymin=0,
			legend cell align=left,
			legend style={
					at={(0.98,0.97)},
					anchor=north east,
					column sep=4pt
			},
			error bars/y dir=both,
			error bars/y explicit,
			legend columns=1,
			scaled y ticks = false,
			ytick={0,100,200,300,400},
			visualization depends on={\thisrow{pores_b_std} \as \offset},
			nodes near coords,
			node near coords style={
				shift={(axis direction cs:0,\offset)},
				rotate=90,
				anchor=west,
				/pgf/number format/.cd,fixed zerofill,precision=1,
			},
			yticklabel pos=right,
		]
			\addplot[
				orange,fill=orange,mark=none,
				error bars/.cd,
				error bar style={black},
			] table [ x expr=\coordindex, y=pores_a_mean, y error=pores_a_std, col sep=comma] {figures/ltdsc/C3S_ice_vol_dev.csv};

			\addplot[
				cyan,fill=cyan,mark=none,
				error bars/.cd,
				error bar style={black}
			] table [ x expr=\coordindex, y=pores_b_mean, y error=pores_b_std, col sep=comma] {figures/ltdsc/C3S_ice_vol_dev.csv};

			% \addplot[
			% 	blue,fill=blue,mark=none,
			% 	error bars/.cd,
			% 	error bar style={black},
			% ] table [ x expr=\coordindex, y=pores_c_mean, y error=pores_c_std, col sep=comma] {figures/ltdsc/C3S_ice_vol_dev.csv};

			\legend{gel pores, capillary pores}%, total pores}
		\end{axis}
	  \end{tikzpicture}
	  \caption{Ice volume of \ac{C3S}}
	  \label{fig:ltdsc_result_C3S_melt_bar}

	\end{subfigure}

    \caption{
		Ice volume per specimen weight for gel and capillary pores of \ac{C3S} and \ac{C3S} from \qtyrange{1}{28}{\day} of hydration.
		The data were calculated from three \ac{LTDSC}-measurements for each value.
		The baseline correction was done using temperature range from a lower range of \qtyrange{-21.5}{-20.5}{\celsius} (capillary pores) and \qtyrange{-57}{-56}{\celsius} (gel pores) and \qtyrange{9}{10}{\celsius} in the upper range.
	}
	\label{fig:ltdsc_result_melt_bar}
\end{figure}


\begin{figure}[htb]
    \centering

	\begin{tikzpicture}
		\begin{axis}[
			ylabel={gel pores (\unit{\micro\liter\per\gram})},
			xlabel={capillary pores (\unit{\micro\liter\per\gram})},
			ymin=0,
			xmin=100,
			width=.8\textwidth,
			height=7cm,
			grid=major,
			enlarge x limits=0.15,
			enlarge y limits=0.15,
			colorbar,
			colorbar style={
				ylabel=hydration time in \unit{\day}
			},
			colormap/viridis,
			point meta min=1,
			point meta max=28,
		]

		% --- PLOT 1: C3S (Kreise) ---
		\addplot [
			scatter, only marks, mark=*, mark size=4pt,
			scatter/use mapped color={fill=mapped color, stroke=mapped color}
		] table [
			y=pores_a_mean, x=pores_b_mean, col sep=comma,
			point meta=\thisrow{a}
		] {figures/ltdsc/C3S_ice_vol_dev.csv};
		\addlegendentry{\ac{C3S}}

		% --- PLOT 2: C2S (Quadrate) ---
		\addplot [
			scatter, only marks, mark=square*, mark size=4pt, % Anderer Marker
			scatter/use mapped color={fill=mapped color, stroke=mapped color}
		] table [
			y=pores_a_mean, x=pores_b_mean, col sep=comma,
			point meta=\thisrow{a}
		] {figures/ltdsc/C2S_ice_vol_dev.csv};
		\addlegendentry{\ac{C2S}}

		\end{axis}
	\end{tikzpicture}

    \caption{
		Correlation between gel and capillary pore volume in \ac{C3S} and \ac{C2S} up to a hydration time of \qty{28}{\day}.
	}
	\label{fig:ltdsc_result_melt_correlation}
\end{figure}



As shown in \zcref{fig:ltdsc_result_C2S_melt_bar}, for \ac{C2S} the ice volume of the gel pores was calculated to be \qty{1.2}{\micro\liter\per\gram} at \qty{7}{\day} up to \qty{18.3}{\micro\liter\per\gram} after \qty{28}{\day} of hydration, while the capillary pores were reduced from \num{290} to \qty{233}{\micro\liter\per\gram} in the same time.
In contrast, \ac{C3S} has a higher ice volume accountable to the gel pores (\zcref{fig:ltdsc_result_C3S_melt_bar}).
The gel pore ice volume was increasing from from \qty{6.6}{\micro\liter\per\gram} at an age of \qty{1}{\day} to \qty{30.3}{\micro\liter\per\gram} after \qty{28}{\day} of hydration.
In comparison to \ac{C2S}, less than half of the ice volume corresponding to capillary pores was detected in \ac{C3S} is after \qty{7}{\day} of hydration.
The capillary pore space was reduced by approximately  \qty{50}{\percent} from \qty{208}{\micro\liter\per\gram} (\qty{1}{\day}) to \qty{109}{\micro\liter\per\gram} (\qty{28}{\day}).
Furthermore, a significant slow down of gel pore development and reduction in capillary pores can be seen in \ac{C3S} after \qty{14}{\day} of hydration.

As expected from the freezing curves, the overall ice volume and the ice volume correlated to the gel pores is lower for \ac{C3S} compared to \ac{C2S}.
Both materials indicate a steady increase in the gel pore and a reduction of capillary ice volume over time.
%As already demonstrated in \zcref{fig:hyd-rim-belite}, the hydrate layer increases in size, filling up the capillary pores.
%At the same time, the growth of \ac{CSH} creates new gel porosity.

These observations fit well to the microstructural development visible in \ac{SEM}-\ac{BSE} images of both materials, where the porosity is reduced, since the hydrate phases grow into them causing new gel porosity (see \zcref{fig:hydration-rim-results}).

\begin{itemize}
	\color{red}
	\item difficult to calculate the pore volume fraction (in contrast to the pore volume per mass)
	\item volume fraction of hydrate phase and unhydrated phase not fully known
	\item can be assumed using Area fraction from SEM images as described in \cite{kleiner.2024a} and \zcref{eq:}.
\end{itemize}

\FloatBarrier
\subsection{\Acl*{MIP} measurements}

\zcref{fig:mip_c2s, fig:mip_c3s} illustrate the \ac{MIP} results for \ac{C2S} and \ac{C3S} after up to \qty{28}{\day} of hydration in comparison to the second freezing curves of the \ac{LTDSC}-measurements.
For these diagrams, the freeze curves were baseline corrected, which was performed by by calculating the linear regression of the data between \numrange{-55}{-50} and \qtyrange{-27}{-23}{\celsius} and subtracting this regression function from the mean data as shown in \zcref{fig:ltdsc_result_c2s_melt}.

\begin{figure}[!htb]
    \centering

	\begin{tikzpicture}
		\begin{semilogxaxis}[
			blue,
			axis line style={black},
			legend style={
				at={(.03,.95)},
				anchor=north west,
				text=black,
			},
			legend cell align={left},
			legend columns=2,
			xmin=1.5, xmax=5000,
			ymin=-.12, ymax=0,
			y dir=reverse,
			width=.85\linewidth,
			height=6cm,
			grid=major,
			xlabel={temperature in \unit{\celsius}},
			xticklabels={\num{-51.2}, \num{-25.3}, \num{-12.4}, \num{-4.6}, \num{-2.0}},
			xtick={2.5,5,10,25,50},
			ylabel={heat flow $\dot{Q}$ in \unit{\joule\per\second\per\gram}},
			scaled y ticks=false,
			scaled x ticks=false,
			axis x line*=top,
			axis y line*=left,
			ytick={0,-0.05,-0.1,-0.15,-0.2},
			yticklabel style={
					/pgf/number format/fixed,
					/pgf/number format/precision=2,
					/pgf/number format/fixed zerofill,
			},
		]

			\addplot[dotted, thick, samples=2, darkgrey,forget plot] coordinates {(3.19,-.2)(3.19,.2)};
			\node[darkgrey, rotate=90, anchor=west] at (axis cs: 2.50,-.035) {\scriptsize <\,\qty{-40}{\celsius}!};

			\addplot [ no marks, orange, forget plot, opacity=0.5] table [each nth point=4, filter discard warning=false, unbounded coords=discard,x index=3, y index=4, col sep=comma] {figures/ltdsc/C2S_7d_mean_freeze_segment.csv};
			\addplot [ no marks, red, forget plot, opacity=0.5] table [each nth point=4, filter discard warning=false, unbounded coords=discard,x index=3, y index=4, col sep=comma] {figures/ltdsc/C2S_14d_mean_freeze_segment.csv};
			\addplot [ no marks, green, forget plot, opacity=0.5] table [each nth point=4, filter discard warning=false, unbounded coords=discard,x index=3, y index=4, col sep=comma] {figures/ltdsc/C2S_21d_mean_freeze_segment.csv};
			\addplot [ no marks, blue, forget plot, opacity=0.5] table [each nth point=4, filter discard warning=false, unbounded coords=discard,x index=3, y index=4, col sep=comma] {figures/ltdsc/C2S_28d_mean_freeze_segment.csv};

			% the color legend is universal for all experiments. The linestyle in addlegendentry however, could not be changed for some reason.
            \addplot [ no marks, orange ]  coordinates {(0.1,0) (.9,0)};
            \addlegendentry{\qty{7}{\day}};
            \addplot [ no marks, red ]  coordinates {(0.1,0) (.9,0)};
            \addlegendentry{\qty{14}{\day}};
            \addplot [ no marks, green ]  coordinates {(0.1,0) (.9,0)};
            \addlegendentry{\qty{21}{\day}};
            \addplot [ no marks, blue ]  coordinates {(0.1,0) (.9,0)};
            \addlegendentry{\qty{28}{\day}};

		\end{semilogxaxis}
		\begin{semilogxaxis}[
			orange,
			axis line style={black},
			legend style={
				at={(.97,.95)},
				anchor=north east,
				text=black,
			},
			legend columns=-1,
			legend cell align={right},
			xmin=1.5, xmax=5000,
			ymin=0, ymax=1.6,
			width=.85\linewidth,
			height=6cm,
			grid=none,
			xlabel={\color{black} pore radius $r_\text{P}$ in \unit{\nano\meter}},
			ylabel style={align=center},
			ylabel={$dV/dlogR$\\pore volume in \unit{\milli\liter\per\gram}},
			grid style={line width=.1pt, draw=gray!10},
			yticklabel style={
					/pgf/number format/fixed,
					/pgf/number format/precision=1,
					/pgf/number format/fixed zerofill,
			},
			xticklabels={1.5,2.5,5,10,25,50, 100, 1000, 10000},
			xtick={1.5,2.5,5,10,25,50, 100, 1000, 10000},
			ytick={},
			scaled y ticks=false,
			scaled x ticks=false,
			axis y line*=right,
			axis x line*=bottom,
			log basis x=10, % seems to be necessary for "restrict x to domain" to work.
		]
			\addplot [ no marks, red, forget plot, dotted] table [skip first n=1, y index=4, x expr={1000*\thisrowno{1}}, col sep=semicolon] {figures/mip/HG C2S 14d.CSV};
			\addplot [ no marks, blue, forget plot, dotted] table [skip first n=1, y index=4, x expr={1000*\thisrowno{1}}, col sep=semicolon] {figures/mip/HG C2S 28d.CSV};


			%the domain restriction is somewhat weird. It should be 10:10000, but that does not work. 1:10000 at least gave the correct result.
			\addplot [ thick, no marks, red, forget plot, restrict x to domain=1:10000] table [skip first n=1, y index=4, x expr={1000*\thisrowno{1}}, col sep=semicolon] {figures/mip/HG C2S 14d.CSV};
			\addplot [ thick, no marks, blue, forget plot, restrict x to domain=1:10000] table [skip first n=1, y index=4, x expr={1000*\thisrowno{1}}, col sep=semicolon] {figures/mip/HG C2S 28d.CSV};

			% distinguish the experiments in the legend
			\addplot [ no marks, black, opacity=0.5] coordinates {(0.1,0) (.9,0)};
			\addlegendentry{\acs*{LTDSC}};
			\addplot [ thick, no marks, black, solid] coordinates {(0.1,0) (.9,0)};
			\addlegendentry{\acs*{MIP}};

		\end{semilogxaxis}
	\end{tikzpicture}
	\caption{{\color{orange}\ac{MIP}-analysis} of \ac{C2S} after \num{14} and \qty{28}{\day} of hydration compared with selected {\color{blue}\ac{LTDSC}-curves}. Note that pore sizes below \qty{3.19}{\nano\meter} are undefined as noted in \zcref{eq:brun_freeze}.}
    \label{fig:mip_c2s}
\end{figure}

\begin{figure}[!htb]
    \centering
	\begin{tikzpicture}
		\begin{semilogxaxis}[
				blue,
				axis line style={black},
				legend style={
					at={(.02,.96)},
					anchor=north west,
					text=black,
				},
				legend cell align={left},
				legend columns=2,
				xmin=1.5, xmax=5000,
				ymin=-.12, ymax=0,
				y dir=reverse,
				width=.85\linewidth,
				height=6cm,
				grid=major,
				xlabel={temperature in \unit{\celsius}},
				xticklabels={\num{-51.2}, \num{-25.3}, \num{-12.4}, \num{-4.6}, \num{-2.0}},
				xtick={2.5,5,10,25,50},
				ylabel={heat flow $\dot{Q}$ in \unit{\joule\per\second\per\gram}},
				scaled y ticks=false,
				scaled x ticks=false,
				axis x line*=top,
				axis y line*=left,
				ytick={0,-0.05,-0.1,-0.15,-0.2},
				yticklabel style={
						/pgf/number format/fixed,
						/pgf/number format/precision=2,
						/pgf/number format/fixed zerofill,
				},
			]

			\addplot[dotted, thick, samples=2, darkgrey,forget plot] coordinates {(3.19,-.2)(3.19,.2)};
			\node[darkgrey, rotate=90, anchor=west] at (axis cs: 2.50,-.020) {\scriptsize <\,\qty{-40}{\celsius}!};

            \addplot [ no marks, yellow, opacity=0.5, forget plot] table [each nth point=4, filter discard warning=false, unbounded coords=discard,x index=3, y index=4, col sep=comma] {figures/ltdsc/C3S_1d_mean_freeze_segment.csv};
            \addplot [ no marks, orange, opacity=0.5, forget plot] table [each nth point=4, filter discard warning=false, unbounded coords=discard,x index=3, y index=4, col sep=comma] {figures/ltdsc/C3S_7d_mean_freeze_segment.csv};
            \addplot [ no marks, red, opacity=0.5, forget plot] table [each nth point=4, filter discard warning=false, unbounded coords=discard,x index=3, y index=4, col sep=comma] {figures/ltdsc/C3S_14d_mean_freeze_segment.csv};
            \addplot [ no marks, green, opacity=0.5, forget plot] table [each nth point=4, filter discard warning=false, unbounded coords=discard,x index=3, y index=4, col sep=comma] {figures/ltdsc/C3S_21d_mean_freeze_segment.csv};
            \addplot [ no marks, blue, opacity=0.5, forget plot] table [each nth point=4, filter discard warning=false, unbounded coords=discard,x index=3, y index=4, col sep=comma] {figures/ltdsc/C3S_28d_mean_freeze_segment.csv};

			% the color legend is universal for all experiments. The linestyle in addlegendentry however, could not be changed for some reason.
            \addplot [ no marks, yellow ]  coordinates {(0.1,0) (.9,0)};
            \addlegendentry{\qty{1}{\day}};
            \addplot [ no marks, orange ]  coordinates {(0.1,0) (.9,0)};
            \addlegendentry{\qty{7}{\day}};
            \addplot [ no marks, red ]  coordinates {(0.1,0) (.9,0)};
            \addlegendentry{\qty{14}{\day}};
            \addplot [ no marks, green ]  coordinates {(0.1,0) (.9,0)};
            \addlegendentry{\qty{21}{\day}};
            \addplot [ no marks, blue ]  coordinates {(0.1,0) (.9,0)};
            \addlegendentry{\qty{28}{\day}};

		\end{semilogxaxis}
		\begin{semilogxaxis}[
			orange,
			axis line style={black},
			legend style={
				at={(.98,.96)},
				anchor=north east,
				text=black,
			},
			legend columns=2,
			legend cell align={right},
			xmin=1.5, xmax=5000,
			ymin=0, ymax=.45,
			width=.85\linewidth,
			height=6cm,
			grid=none,
			xlabel={pore radius $r_\text{P}$ \unit{\nano\meter}},
			ylabel style={align=center},
			ylabel={$dV/dlogR$\\pore volume in \unit{\milli\liter\per\gram}},
			grid style={line width=.1pt, draw=gray!10},
			yticklabel style={
					/pgf/number format/fixed,
					/pgf/number format/precision=1,
					/pgf/number format/fixed zerofill,
			},
			xticklabels={1,2.5,5,10,25,50, 100, 1000, 10000},
			xtick={1,2.5,5,10,25,50, 100, 1000, 10000},
			ytick={},
			scaled y ticks=false,
			scaled x ticks=false,
			axis y line*=right,
			axis x line*=bottom,
			log basis x=10, % seems to be necessary for "restrict x to domain" to work.
		]
			\addplot [ orange, dotted, forget plot] table [skip first n=1, y index=4, x expr={1000*\thisrowno{1}}, col sep=semicolon] {figures/mip/HG C3S 7d.CSV};
			\addplot [ no marks, blue, dotted, forget plot] table [ skip first n=1, y index=4, x expr={1000*\thisrowno{1}}, col sep=semicolon] {figures/mip/HG C3S 28d.CSV};

			%the domain restriction is somewhat weird. It should be 10:10000, but that does not work. 1:10000 at least gave the correct result.
			\addplot [ thick, no marks, orange, forget plot, restrict x to domain=1:10000] table [skip first n=1, y index=4, x expr={1000*\thisrowno{1}}, col sep=semicolon] {figures/mip/HG C3S 7d.CSV};
			\addplot [ thick, no marks, blue,forget plot, restrict x to domain=1:10000] table [skip first n=1, y index=4, x expr={1000*\thisrowno{1}}, col sep=semicolon] {figures/mip/HG C3S 28d.CSV};

			% distinguish the experiments in the legend
			\addplot [ no marks, black, opacity=0.5] coordinates {(0.1,0) (.9,0)};
			\addlegendentry{\acs{LTDSC}};
			\addplot [ thick, no marks, black, solid] coordinates {(0.1,0) (.9,0)};
			\addlegendentry{\acs{MIP}};

		\end{semilogxaxis}
	\end{tikzpicture}
	\caption{
		{\color{orange}\ac{MIP}-analysis} of \ac{C3S} after \num{7} and \qty{28}{\day} of hydration compared with selected {\color{blue}\ac{LTDSC}-curves}.
		Note that pore sizes below \qty{3.19}{\nano\meter} are undefined as noted in \zcref{eq:brun_freeze}.
	}
    \label{fig:mip_c3s}
\end{figure}

The \ac{LTDSC} curves above \qty{-12.4}{\celsius} ($r_P$=\qty{10}{\nano\meter}) are not reliable due to system inertia.
This is because the rapid freezing rate shifts the capillary pore space and larger pores become indistinguishable.
However, \ac{MIP} is considered fairly reliable above \qty{10}{\nano\meter}, while smaller pores are hardly detectable.
Therefore, both methods cover different pore space ranges.
\ac{LTDSC} indicates the development of the gel pores, while \ac{MIP} curves show the development of the capillary and larger pores.

The \ac{MIP} curves show a significant reduction in the pore size distribution for both materials from the first reliable measurable age (\ac{C2S}: \qty{14}{\day}; \ac{C3S}: \qty{7}{\day}) up to \qty{28}{\day} of hydration.
As already indicated by \ac{LTDSC}, the capillary pore space of \ac{C2S} is significantly higher than that of \ac{C3S} (not the scale difference of both diagrams).

%\section{Discussion}
%\color{red}
%TODO: compare results of the methods
%\color{black}

\section{Conclusion}

The characterisation of porosity in cementitious materials necessitates a multi-method approach to cover the full range of pore sizes.
Benchmarking against porous glass standards revealed that \ac{SEM}-based pore size analysis is highly sensitive to segmentation parameters.
In contrast, \ac{LTDSC} offers reliable data for the overall porosity, which can also be distinguished between gel and capillary pores.
However, the assignment to an actual pore size using the approximation by \textcite{brun.1977} seems to be unsuitable, since the water in the gel pores freezes significantly below the boundary temperature \qty{-40}{\degreeCelsius} of the equation.

In terms of microstructural imaging, \ac{ArBIB} sectioning was found to be superior to mechanical polishing.
It revealed the solid, star-shaped cross-section of \ac{CSH} needles and preserved the native pore structure without requiring resin embedding or cryo-cooling.
Regarding the hydration of \ac{C3S} and \ac{C2S}, both \ac{LTDSC} and \ac{MIP} measurements consistently indicated a densification of the microstructure over time, characterised by an increase in gel porosity and a reduction in capillary porosity.
In accordance with the results based on image analysis shown in \zcref{chap:hydrate-rim}, \ac{C3S} showed a significantly higher volume of gel pores and a lower volume of capillary pores compared to \ac{C2S} at the same hydration age, reflecting its faster reaction kinetics.

However, none of the methods reliably covers the entire pore range.
Furthermore, the results of the methods examined are hardly comparable with each other.